\section{Conclusion}

Currently, low-level code is almost exclusively written in C.  Pit is a much friendlier language for writing low-level
code, and there is no penalty when the programmer chooses to use only primitive variables.  When the programmer chooses
to use auto variables, there is a cost (we don't know how much yet), but this cost is mitigated in two ways.  First, a
decent optimizer will eliminate a substantial amount of the overhead, and second, the programmer may choose to use
primitive variables only where the benefit is worthwhile.  Because of the relation of Zipf's Law to execution time, the
penalty will not be important in most of the code.  When efficiency is necessary, Pit provides an exceptionally easy
way to interface auto variables with primitive variables.  When using auto variables, the code is completely protected
from buffer overruns, integer overflow bugs, and other memory allocation bugs.  Object files generated from modules
written in Pit are far easier to move between operating systems, as recompiling is usually not required.  These
considerations make Pit a superior language to C for low-level programming.
